\section{Model}

We propose \ImitatorEOne{}, an autoregressive architecture that predicts exact
\Gemma{} token IDs from 2D keypoints, extending the modular premise of our
earlier embedding-based \system{} \citep{velasquez2026imitator} with the
discrete token-ID objective introduced in Section~3.

\paragraph{Visual encoder.}
Three ST-GCN blocks process each keypoint graph, followed by a 128-channel
projection, node averaging, two temporal convolutions, and a one-layer, four-head
Transformer.  Fixed sinusoidal position is added before self-attention
\citep{vaswani2017attention}.  Without it, segment permutation left 86.3\% of
predictions unchanged; with it, content can be bound to video position.  ST-GCN
and the spatial projection use signer-clean initialization and remain frozen;
temporal layers unfreeze after five epochs.

\paragraph{Token decoder.}
A two-layer, four-head causal Transformer (hidden size 128; feed-forward size 512)
attends to visual memory.  Input and output embeddings are tied.  Greedy decoding
starts at BOS, emits at most 32 content tokens, and stops at EOS.  The task contains
118 \Gemma{} content IDs; with BOS, EOS, and PAD, a bijection gives a 121-row
softmax.  Rows map back to the original IDs, so this is computational
reparameterization rather than a new sign vocabulary.

\paragraph{Language stage.}
The same tokenizer decodes emitted IDs for a fixed \Gemma{} 3n E2B instruction
prompt \citep{gemma3n}.  We do not jointly fine-tune this downstream model; this
is an experimental-design choice.  We report transcription before the LLM and
correction afterward, preventing fluent output from masking visual error.
